\documentclass[11pt]{article}
\usepackage[margin=1in]{geometry}
\usepackage[T1]{fontenc}
\usepackage{lmodern}
\usepackage{microtype}
\usepackage{amsmath,amssymb,amsthm,mathtools}
\usepackage{enumitem}
\usepackage[hidelinks]{hyperref}
\usepackage{xcolor}
\usepackage{booktabs}

\newtheorem{proposition}{Proposition}
\newtheorem{lemma}{Lemma}
\newtheorem{corollary}{Corollary}
\newtheorem{remark}{Remark}
\newcommand{\ord}{\operatorname{ord}}
\newcommand{\Spec}{\operatorname{Spec}}
\newcommand{\Res}{\operatorname{Res}}
\newcommand{\A}{\mathbb A}
\newcommand{\Q}{\mathbb Q}
\newcommand{\F}{\mathbb F}
\newcommand{\kbar}{\overline{k}}

\title{Corrections and Proof Repairs for\\
\emph{Boundary Belyi Covers and Normal-Jet Obstructions in the Plane Jacobian Problem}, v13}
\author{Audit repair note}
\date{30 July 2026}

\begin{document}
\maketitle

\begin{abstract}
This note records source-level repairs to the v13 working manuscript.  It
corrects the codomain in Proposition~3.3, replaces the proof of the
contact-degree formula with an invariant wedge-order argument, supplies a
non-computational proof of the asserted $A_{21}$ monodromy, states the exact
logical requirement behind the Macaulay minor in Theorem~C.6, and replaces
the finite-etale sentence in Theorem~C.8 by a finite-plus-Nakayama argument.
The degree-$17$ secondary Belyi cover survives the contact-degree repair:
in the stored $(8,28)$ chart the computed positive wedge order itself forces
the residual target parameter to be constant.  This note does not replay the
large certificate archive and does not remove the upstream exhaustiveness
condition needed for a global degree-below-$125$ theorem.
\end{abstract}

\section{Status and scope}

The v13 PDF is a working manuscript whose title page is dated 22 July 2026;
29 July 2026 is the public release date appearing in the filename and site
manifest.  These dates should be displayed separately.

The repairs below distinguish three levels of assertion:
\begin{enumerate}[label=(\roman*)]
  \item hand proofs that can be repaired directly;
  \item exact claims about named finite systems whose certificate semantics
        must be stated completely; and
  \item the global degree-below-$125$ implication, which remains conditional
        on an audited exhaustive Newton reduction and attachment argument.
\end{enumerate}

\section{Corrected codomain for Proposition 3.3}

Let $A=\deg p$ and $B=\deg q$, and write the reduced face equation as
\[
 F(p,q)=Npq-m u\,p q'+n u\,p'q=\text{constant}.
\]
Variations fixing the constant coefficients lie in
\[
 U:=u k[u]_{\le A}\oplus u k[u]_{\le B}.
\]
The differential is
\[
 L(\alpha,\beta)=N(\alpha q+p\beta)
 -m u(\alpha q'+p\beta')+n u(\alpha' q+p'\beta).
\]
The constant coefficient of $L(\alpha,\beta)$ vanishes.  Its coefficient at
$u^{A+B}$ is
\[
 (N-mB+nA)(\alpha_Aq_B+p_A\beta_B),
\]
and the endpoint degree relation for the face equation is
$N-mB+nA=0$.  Hence
\[
 L:U\longrightarrow W:=u k[u]_{\le A+B-1}.
\]

\begin{proposition}[Replacement for Proposition 3.3]
For each explicit reduced cover in Section~3, the map
\[
 L:u k[u]_{\le A}\oplus u k[u]_{\le B}
   \longrightarrow u k[u]_{\le A+B-1}
\]
is surjective and its kernel is the source-scaling line spanned by
$(up',uq')$.  Consequently the normalized quotient solution is a reduced
isolated point after quotienting by source scaling.
\end{proposition}

\noindent
The rank table in v13 proves this corrected statement: the domain has
dimension $A+B$, the displayed rank is $A+B-1=\dim W$, and the nonzero
source-scaling vector already lies in the kernel.

\section{Contact degree: corrected statement and proof}

Let $D=(t=0)$ be a smooth boundary component with a parameter $s$ along
$D$.  For a meromorphic two-form, $\ord_D$ denotes the $t$-order of its
coefficient in $dt\wedge ds$.

\begin{lemma}[Leading wedge order and the Section~5 resonances]
Suppose
\[
 \pi=t^a c(s)+O(t^{a+1}),\qquad
 \tau=\tau_0(s)+\sum_{j\ge1}t^j h_j(s),
 \qquad c\ne0.
\]
If $\ord_D(d\pi\wedge d\tau)>a-1$, then $d\tau_0=0$ in $k(D)$.
In the normalized Section~5 chart, where $c(s)=s/(s-1)$, a rational lower
coefficient invisible to the wedge occurs exactly at an order $j=ar$ and is
removed by the target shear $\tau\mapsto\tau-C\pi^r$.  After all such lower
resonances have been removed, if $e$ is the first remaining normal order,
then
\[
 \ord_D(d\pi\wedge d\tau)=a+e-1.
\]
\end{lemma}

\begin{proof}
The coefficient of $t^{a-1}dt\wedge ds$ in $d\pi\wedge d\tau$ is
$a c(s)\tau_0'(s)$.  Thus an order greater than $a-1$ forces
$\tau_0'=0$.  At order $a+j-1$, the leading coefficient contributed by
$t^jh_j$ is
\[
 a c h_j'-j c'h_j.
\]
It vanishes exactly when $h_j=Cc^{j/a}$.  For $c=s/(s-1)$, the divisor of
$c$ is $(0)-(1)$, so this expression is rational exactly when $a\mid j$.
Writing $j=ar$, it is the leading term of the target shear $C\pi^r$ and may
be subtracted without changing $d\pi\wedge d\tau$.  Once all lower
resonant terms have been removed, the first remaining order $e$ has nonzero
coefficient, giving $a+e-1$.
\end{proof}

\begin{proposition}[Corrected contact-degree formula]
Suppose coprime positive integers $m,n$ satisfy
\[
 v_E(P)=-m,\quad v_E(Q)=-n,\quad
 v_D(P)=-ma,\quad v_D(Q)=-na.
\]
Choose $c,d\in\mathbb Z$ with $dn-cm=1$ and set
\[
 \pi=P^c/Q^d,\qquad \tau=Q^m/P^n.
\]
Assume $x=t^{-1}$ up to a unit,
\[
 dP\wedge dQ=x^\kappa\,dx\wedge dy
\]
up to a unit along $D$, and use the normalized Section~5 coordinate
$\pi=t^a s/(s-1)+O(t^{a+1})$.  Put
$K_D=\ord_D(dx\wedge dy)$ and
\[
 e_*:=a(m+n)-\kappa+K_D+1.
\]
If $e_*>0$, then $\tau|_D$ is constant.  After removing all lower resonant
target shears, the first non-shear contact degree is
\[
 \boxed{e=e_*.}
\]
In the regular toric chart used in v13, $K_D=-2$, so
\[
 \boxed{e=a(m+n)-\kappa-1.}
\]
\end{proposition}

\begin{proof}
Direct differentiation gives
\[
 d\pi\wedge d\tau
 =-P^{c-n-1}Q^{m-d-1}\,dP\wedge dQ.
\]
The monomial factor has $D$-order $a(m+n+1)$, because
$dn-cm=1$.  Therefore
\[
 N:=\ord_D(d\pi\wedge d\tau)
   =a(m+n+1)-\kappa+K_D.
\]
The inequality $e_*>0$ is equivalent to $N>a-1$.  The preceding lemma
therefore forces $\tau|_D$ to be constant.  After resonant target shears are
removed, the same lemma gives $N=a+e-1$, hence
$e=N-a+1=e_*$.
\end{proof}

\begin{corollary}[The stored $(4,17)$ contact survives]
For the stored chart with $(m,n,a,\kappa,K_D)=(2,3,4,2,-2)$,
\[
 e=4(2+3)-2-1=17>0.
\]
Thus the wedge order itself forces constant residual $\tau$, and the
shear-normalized transport theorem applies with $(a,e)=(4,17)$.  The
resulting degree-$17$ Belyi map and passport
\[
 (4^4,1),\qquad (17),\qquad (5,1^{12})
\]
are unaffected by this correction.
\end{corollary}

\section{A non-computational proof of the $A_{21}$ monodromy}

Let $G$ be the monodromy group of a connected degree-$21$ cover with
passport
\[
 (2^{10},1),\qquad (3^7),\qquad (17,1^4).
\]
The following replaces the permutation-certificate portion of
Proposition~C.1; the exact count of five covers and their arithmetic orbit
remain computational statements.

\begin{proposition}
The monodromy group is $A_{21}$.
\end{proposition}

\begin{proof}
The three branch permutations are even, so $G\subseteq A_{21}$.  The cover
is connected, hence $G$ is transitive.  A nontrivial block in degree $21$
would have size $3$ or $7$.  The $17$-cycle induces the identity on the set
of blocks, because neither $S_7$ nor $S_3$ contains an element of order
$17$.  It would therefore preserve every block setwise.  This is impossible:
a preserved block is a union of orbits of the $17$-cycle, but its unique
nontrivial orbit has size $17$, larger than either possible block.  Thus $G$
is primitive.  Jordan's theorem now applies: a primitive subgroup of
$S_{21}$ containing a prime $17$-cycle, with $17\le21-3$, contains
$A_{21}$.  Hence $G=A_{21}$.
\end{proof}

\section{Theorem C.6: exact certificate semantics}

A nonzero square minor proves a rank statement.  It proves a unit-ideal
statement only after the domain and the \emph{complete} target basis of the
Macaulay map have been specified.

Let the six selected polynomials be $F_1,\dots,F_6$.  The certificate should
define finite multiplier spaces $V_i$ and a target monomial space $W$ such
that $F_iV_i\subseteq W$, together with the map
\[
 \mu:\bigoplus_{i=1}^6V_i\longrightarrow W,
 \qquad (H_i)_i\longmapsto\sum_i H_iF_i.
\]
The logically sufficient replacement text is:

\begin{quote}
The archive stores the full ordered basis of $W$, consisting of exactly
$7121$ monomials and including the constant monomial $1$, and the ordered
basis of $\bigoplus_iV_i$, consisting of $10824$ columns.  The certified
$7121\times7121$ minor uses every row of $W$ and the listed pivot columns.
Its reduction has determinant $859\pmod{2053}$, so the minor is nonzero in
characteristic zero.  Hence $\mu$ is surjective; in particular
$1\in\operatorname{im}\mu\subseteq(F_1,\dots,F_6)$.
\end{quote}

\noindent
This wording is valid only if the archive replay confirms the phrases
``full ordered basis,'' ``including $1$,'' and ``uses every row.''  If the
$7121$ rows are merely a subset of a larger target basis, the minor alone
does not prove $1\in(F_1,\dots,F_6)$; an explicit solution of
$\mu(H_i)=1$ or a Nullstellensatz cofactor identity is then required.
The explicit cofactor identities stated in Theorem~C.7 are already a
complete proof interface for the two stored branch ideals, subject to the
separate correctness of the reduction to those branches.

\section{Theorem C.8: finite plus Nakayama, not unproved flatness}

Let $R$ be the unramified local DVR obtained by localizing the coefficient
ring at $(2053,u-216)$, with residue field $k$.  Let $\mathcal T_R$ be the
proper toric model used by the certificate, and put
\[
 Z:=V(g_1,\dots,g_5)\subseteq\mathcal T_R.
\]

\begin{proposition}[Replacement for Theorem C.8]
Assume the face-initial-ideal replay shows that the special fiber $Z_k$ has
no toric-boundary points and is the reduced finite scheme of length $296$
represented by the archived multiplication matrices.  If multiplication by
$\rho$ on $\Gamma(Z_k,\mathcal O_{Z_k})$ has nonzero determinant, then
\[
 V_{\overline{K_0}}(\rho,g_1,g_2,g_3,g_4,g_5)=\varnothing.
\]
For the stored specialization the determinant is
$682\not\equiv0\pmod{2053}$, so the conclusion follows.
\end{proposition}

\begin{proof}
The morphism $Z\to\Spec R$ is proper.  It is quasi-finite on the closed
fiber.  The non-quasi-finite locus is closed in $Z$, and its image is closed
because the morphism is proper.  That image cannot be the generic point
alone in the spectrum of a local DVR, and it does not contain the closed
point.  Hence it is empty.  Thus $Z\to\Spec R$ is proper and quasi-finite,
therefore finite.

Write $A=\Gamma(Z,\mathcal O_Z)$, a finite $R$-algebra.  The determinant
condition says that multiplication by $\bar\rho$ on $A/\mathfrak mA$ is an
automorphism.  Consequently
\[
 (A/\rho A)\otimes_R k=0.
\]
The finite $R$-module $A/\rho A$ is therefore zero by Nakayama's lemma.
Thus $\rho$ is a unit in $A$, so the six equations have no common point on
$Z$ over $R$, on its generic fiber, or after any field extension.  In
particular the common zero locus over $\overline{K_0}$ is empty.
\end{proof}

\begin{remark}
Finite flatness or finite etaleness may also be provable from the smooth
torus complete-intersection structure and a Jacobian-rank check, but neither
property is needed for the no-common-zero conclusion.  The replacement
proof uses only properness, finiteness of the special fiber, and invertibility
of multiplication by $\rho$ modulo the maximal ideal.
\end{remark}

\section{Required public scope language}

The manuscript explicitly states that it does not supply a new independent
proof of the degree-$125$ lower bound.  The corresponding site collection
should therefore use the following statement:

\begin{quote}
For each displayed terminal system, the exact certificates establish the
claimed emptiness.  A global conclusion that every plane Keller map of
maximum coordinate degree below $125$ is an automorphism additionally
requires an independently replayed and audited proof that the upstream
Newton reduction, saturations, normalizations, branch eliminations, and
chart attachments exhaust every candidate.
\end{quote}

\noindent
The degree-below-$125$ statement may be reported only as an external,
source-attributed announcement or literature claim, with its independent
review status shown separately.  It must not be presented as a theorem proved
by Program~6 or by the v13 manuscript.

\section{Disposition of the audited items}

\begin{center}
\begin{tabular}{@{}p{0.23\textwidth}p{0.67\textwidth}@{}}
\toprule
Item & Repair status \\
\midrule
Proposition 3.3 & Repaired by specifying the codomain
$u k[u]_{\le A+B-1}$. \\
Proposition 5.2 & Repaired.  Positivity of the computed wedge order forces
constant residue; the $(4,17)$ application survives. \\
Proposition C.1 & $A_{21}$ now follows from transitivity, primitivity,
parity, and Jordan's theorem; no permutation certificate is needed. \\
Theorem C.6 & Logical implication repaired, but the archive must explicitly
confirm that the $7121$ rows are the complete target basis containing $1$. \\
Theorem C.8 & Repaired by properness, finiteness, and Nakayama; geometric
emptiness is stated over $\overline{K_0}$. \\
Degree below $125$ & Remains conditional within Program~6 on the upstream
exhaustiveness and attachment audit. \\
\bottomrule
\end{tabular}
\end{center}

\end{document}
